\documentclass[11pt,a4paper]{article}
\usepackage{amsmath,amssymb,amsthm}
\usepackage{geometry}
\usepackage{hyperref}
\usepackage{booktabs}
\geometry{margin=2.5cm}

\newtheorem{theorem}{Theorem}[section]
\newtheorem{lemma}[theorem]{Lemma}
\newtheorem{corollary}[theorem]{Corollary}
\newtheorem{definition}[theorem]{Definition}
\newtheorem{remark}[theorem]{Remark}

\title{The Cascade Series --- Part 0\\[0.5em]
\large Scale Variance from Orthogonality:\\
How the Unit Ball Generates $10^{120}$ Orders of Magnitude}
\author{RTAC}
\date{March 2026}

\begin{document}
\maketitle

\begin{abstract}
The unit ball in $\mathbb{R}^d$ has a slicing recurrence
$V_{d+1}/V_d = \sqrt{\pi}\cdot R(d)$ containing one constant,
$\sqrt{\pi} = \Gamma(1/2)$, forced by orthogonality. The sphere-area decay
rate $p(d)$ has a unique natural zero at $d_0 = 7$. This zero generates
two threshold values $c_1 = \tfrac{1}{2}\ln\pi$ and $c_2 = \sqrt{\pi}$ ---
the two canonical values of the recurrence's unique constant, one in each
of its two canonical forms. The orbit terminates because a first-order
recurrence has exactly two canonical forms and one bridge map between them.
The thresholds select $d_1 = 19$ and $d_2 = 217$.

The Gamma function produces exactly four distinguished dimensions in the
cascade: the volume maximum ($d_V = 5$), the sphere-area maximum
($d_0 = 7$), and the two thresholds ($d_1 = 19$, $d_2 = 217$). No fifth
exists.

The cascade invariant
$I_0 = \Omega_{19}\times\Omega_{217} = 1.2051\times 10^{-120}$
is the unique product of the threshold sphere areas under three deductive
constraints from the recurrence's structure. The formula for the hierarchy
is $\pi e^{2\sqrt{\pi}}(2\sqrt{\pi}-1)/\ln 10\approx 120.3$.

Every step is a theorem about the Gamma function. No physics enters.
\end{abstract}

\section{The Question}

How much scale variance can pure mathematics produce from nothing?

The unit ball $B^d = \{x\in\mathbb{R}^d : |x|\leq 1\}$ exists in every
dimension. Its boundary sphere has area
$\Omega_{d-1} = 2\pi^{d/2}/\Gamma(d/2)$. At $d = 7$:
$\Omega_7\approx 33$. At $d = 217$:
$\Omega_{217}\approx 10^{-120}$. The ratio exceeds $10^{121}$. The unit
ball generates 121 orders of magnitude of scale variance with no free
parameters and no input beyond the definition of orthogonal dimensions.

This paper extracts two distinguished scales from this variance and proves
their product is unique.

\section{The Recurrence}

Cut $B^{d+1}$ perpendicular to one axis. Each cross-section at height $x$
is a $d$-ball of radius $\sqrt{1-x^2}$:
\begin{equation}\label{eq:slice}
V_{d+1} = V_d\int_{-1}^{1}(1-x^2)^{d/2}\,dx.
\end{equation}
Substitute $x = \cos\theta$. The integral becomes
$2\int_0^{\pi/2}\sin^d\theta\,d\theta$: one quarter turn from axis to
equator. Via the Beta function:
\begin{equation}\label{eq:recurrence}
V_{d+1} = V_d\cdot\sqrt{\pi}\cdot R(d), \qquad
R(d) := \frac{\Gamma((d+1)/2)}{\Gamma((d+2)/2)}.
\end{equation}

\begin{theorem}[Orthogonal compression constant]\label{thm:sqrtpi}
$\Gamma(1/2) = \sqrt{\pi}$ is the unique dimension-independent constant in
the recurrence, forced by orthogonality. It appears with coefficient
exactly~1.
\end{theorem}

\begin{proof}
The first argument of $B$ is $\frac{1}{2}$, forced by orthogonality: the
angle between axis and equator is $\pi/2$; the substitution
$x = \cos\theta$ produces the half-power $t^{-1/2}$ in $B(1/2,\cdot)$;
$\Gamma(1/2) = \sqrt{\pi}$. The chain: dimension $\to$ orthogonal $\to$
$\pi/2$ $\to$ $B(1/2,\cdot)$ $\to$ $\sqrt{\pi}$. Every arrow forced.
\end{proof}

\section{Sphere Areas Are Primary}

\begin{theorem}[Boundary dominance]\label{thm:boundary}
$\Omega_{d-1}/V_d = d$ for all $d\geq 1$.
\end{theorem}

\begin{proof}
$\Omega_{d-1} = 2\pi^{d/2}/\Gamma(d/2)$,\;
$V_d = \pi^{d/2}/[(d/2)\Gamma(d/2)]$.\; Ratio: $d$.
\end{proof}

\begin{corollary}\label{cor:primary}
$V_d = \Omega_{d-1}/d$. Volumes are derived from sphere areas. The
recurrence is a first-order multiplicative map between sphere areas:
\[
\Omega_d = \Omega_{d-1}\cdot\tfrac{d+1}{d}\cdot\sqrt{\pi}\cdot R(d).
\]
Sphere area is the unique independent cascade quantity at each level.
Every other cascade object---$V_d$, $R(d)$,
$R_{\mathrm{eff}}(d) = 1/\sqrt{d+3}$---is derived from sphere areas.
Any identification mapping cascade content to energy must therefore
assign energy $C\cdot\Omega_{d-1}$ at each level, where $C$ is a
universal constant: no independent geometric quantity exists to make
$C$ depend on $d$.
\end{corollary}

\section{Irreducibility}

\begin{theorem}\label{thm:irred}
$\sqrt{\pi}$ is the irreducible residue of the recurrence.
\end{theorem}

\begin{proof}
The second difference
$\delta(d) := \ln(\Omega_{d+1}/\Omega_d) - \ln(\Omega_d/\Omega_{d-1})
= \ln(d/(d+1))$. All $\pi$-dependent and $\Gamma$-dependent factors
cancel exactly. $\sqrt{\pi}$ is the only nontrivial constant the
recurrence contains.
\end{proof}

\section{The Natural Zero}

The sphere-area decay rate decomposes as:
\[
p(d) = -\tfrac{1}{2}\ln\pi +
\tfrac{1}{2}\psi\!\left(\tfrac{d+1}{2}\right),
\]
where $\psi$ is the digamma function. Since $\psi^{(1)} > 0$, $p$ is
strictly monotone increasing with a unique zero.

\begin{definition}[Natural zero]
$d_0$ is the unique solution to $p(d_0) = 0$. Stirling:
$d_0\approx 2\pi$; integer $d_0 = 7$. This is the sphere-area maximum.
\end{definition}

\section{The Threshold Pair}

\subsection{Operation A: the zero-crossing forces $c_1$}

\begin{theorem}[Operation A]\label{thm:opA}
The zero-crossing forces $c_1 := \tfrac{1}{2}\ln\pi$.
\end{theorem}

\begin{proof}
$p(d_0) = 0$ gives $\frac{1}{2}\psi((d_0+1)/2) = \frac{1}{2}\ln\pi$.
The right-hand side is the constant part of $p$, forced by
Theorem~\ref{thm:sqrtpi}.
\end{proof}

\begin{theorem}[Uniqueness of $c_1$]\label{thm:c1unique}
Among all scalars derivable from $d_0$, $c_1 = \tfrac{1}{2}\ln\pi$ is the
unique scalar satisfying (a) $c > 0$ (so $p(d) = c$ has a solution beyond
$d_0$) and (b) primitivity ($c$ is a recurrence constant).
\end{theorem}

\begin{proof}
We check all four information classes available from the natural zero.

\emph{Class~1: Values of $p$ and its decomposition at $d_0$.} The zero
gives three scalars: $p(d_0) = 0$, the $d$-dependent part
$q(d_0) = \frac{1}{2}\ln\pi$, and the constant part
$-\frac{1}{2}\ln\pi$. (i)~$p(d_0) = 0$ fails~(a): $p(d) = 0$ has its
solution at $d_0$ itself. (ii)~$q(d_0) = \frac{1}{2}\ln\pi$ satisfies~(a)
since $\frac{1}{2}\ln\pi > 0$, and satisfies~(b) since
$\frac{1}{2}\ln\pi$ is the log-space image of $\sqrt{\pi}$, the
recurrence's unique primitive. This is $c_1$. (iii)~$-\frac{1}{2}\ln\pi < 0$
fails~(a).

\emph{Class~2: Higher derivatives $p^{(n)}(d_0)$.} All fail~(b): they are
not recurrence primitives.

\emph{Class~3: The dimension $d_0$ itself.} $p(d) = d_0$ equates a
dimensionless rate to a dimension count---a type mismatch.

\emph{Class~4: Analytic combinations.} Any combination either reduces to
$q(d_0)$ or introduces quantities that are not recurrence primitives.

Across all four classes, $c_1 = \frac{1}{2}\ln\pi$ is unique.
\end{proof}

\subsection{Two canonical values of one constant}

The recurrence contains one irreducible constant
(Theorem~\ref{thm:irred}). That constant has two canonical
representations, because the recurrence has two canonical forms:

\begin{center}
\begin{tabular}{lll}
\toprule
Form & Expression & Constant \\
\midrule
Multiplicative & $V_{d+1}/V_d = \sqrt{\pi}\cdot R(d)$ &
$c_2 = \sqrt{\pi}$ \\
Logarithmic & $\Delta\ln V_d = \frac{1}{2}\ln\pi + \ln R(d)$ &
$c_1 = \frac{1}{2}\ln\pi$ \\
\bottomrule
\end{tabular}
\end{center}

These are the same constant in two algebraic spaces, related by the
exponential map: $e^{c_1} = e^{\frac{1}{2}\ln\pi} = \sqrt{\pi} = c_2$.

\begin{definition}[Canonical form]\label{def:canonical}
A canonical form of a first-order recurrence is an algebraically equivalent
expression that isolates the dimension-independent content as a single
constant with coefficient~1. For a multiplicative recurrence
$a_{n+1} = K\cdot f(n)\cdot a_n$, the two canonical forms are the
recurrence itself (constant $K$) and its logarithm (constant $\ln K$).
Any other algebraically equivalent expression either has a constant that
is a power, root, or inverse of one of these (hence derived, not
primitive) or fails to isolate the constant with coefficient~1.
\end{definition}

\begin{theorem}[Two canonical values]\label{thm:twovalues}
The recurrence's dimension-independent content is exactly the pair
$(c_1, c_2) = (\tfrac{1}{2}\ln\pi,\;\sqrt{\pi})$: one constant in each
canonical form. No third canonical value exists.
\end{theorem}

\begin{proof}
By Definition~\ref{def:canonical}, a first-order multiplicative recurrence
has exactly two canonical forms. The multiplicative form isolates
$\sqrt{\pi}$; the logarithmic form isolates $\frac{1}{2}\ln\pi$. Any
other rewriting---squaring ($\pi$), inverting ($1/\sqrt{\pi}$), taking
roots ($\pi^{1/4}$)---produces a derived constant, not a primitive one:
the derived constant is a function of $\sqrt{\pi}$ and cannot appear with
coefficient~1 in any first-order form of the recurrence. For the slicing
recurrence: $K = \sqrt{\pi}$, $\ln K = \frac{1}{2}\ln\pi$. These exhaust
the dimension-independent content.
\end{proof}

\begin{remark}
The first threshold comes from applying $p(d) = c$ to the logarithmic
constant $c_1 = \frac{1}{2}\ln\pi$. The second comes from applying the
same construction to the multiplicative constant $c_2 = \sqrt{\pi}$.
Using only one would ignore half the recurrence's constant content. The
bridge between them is not an arbitrary operation: it is the exponential
map that defines the relationship between a multiplicative recurrence and
its logarithm.
\end{remark}

\subsection{Orbit termination}

\begin{theorem}[Orbit termination]\label{thm:terminate}
The orbit of $\{0\}$ under the threshold construction is
$\{0, c_1, c_2\}$. No further element is reachable.
\end{theorem}

\begin{proof}
The orbit is generated by two operations: (A)~reading $c_1$ from the
natural zero, and (B)~passing to the other canonical form via $\exp$.
A first-order multiplicative recurrence has exactly two canonical forms
(Theorem~\ref{thm:twovalues}), connected by exactly one bridge map
($\exp/\ln$). The orbit visits the zero (seed), the logarithmic constant
($c_1$), and the multiplicative constant ($c_2$). No fourth station
exists because there is no third canonical form.

Explicitly: (i)~$\exp(c_2) = e^{\sqrt{\pi}}\approx 5.83$ is not a
recurrence primitive by Theorem~\ref{thm:irred}---the recurrence's
only nontrivial constant is $\sqrt{\pi}$, and $e^{\sqrt{\pi}}$ does not
appear. (ii)~$\ln(c_1) = \ln(\frac{1}{2}\ln\pi) \approx -0.558$ is
negative and fails to generate a threshold beyond $d_0$.
(iii)~Operation~A cannot be reapplied: $p$ has exactly one zero (strict
monotonicity), so there is no second zero at which to read off a new
value.

The orbit terminates at three elements because the recurrence has two
canonical values and one structural zero. This is a property of the
recurrence's order, not a coincidence about specific numbers.
\end{proof}

\subsection{Two thresholds and no more}

\begin{theorem}[Two thresholds]\label{thm:thresholds}
$p(d) = c_1$ has a unique solution at $d_1 = 19$.\;
$p(d) = c_2$ has a unique solution at $d_2 = 217$.
\end{theorem}

\begin{proof}
Strict monotonicity of $p$ gives uniqueness. Discrete verification:
$p(19) = 0.5535 < c_1 = 0.5724 < p(20) = 0.6013$.
$p(217) = 1.77101 < c_2 = 1.77245 < p(218) = 1.77331$.
\end{proof}

\begin{theorem}[No third threshold]
The orbit terminates at $(c_1, c_2)$. Each positive canonical value
determines exactly one threshold dimension. The cascade has exactly two
thresholds.
\end{theorem}

\section{Four Distinguished Dimensions}

\begin{theorem}[Tower completeness]\label{thm:tower}
The Gamma function produces exactly four distinguished dimensions in the
cascade. No fifth exists.
\end{theorem}

\begin{proof}
Two maxima: $V_d$ has a unique maximum at $d_V = 5$ and $\Omega_d$ has a
unique maximum at $d_0 = 7$, each by unimodality. Two thresholds: the
complete constant content generates exactly $(c_1, c_2)$ and the orbit
terminates (Theorem~\ref{thm:terminate}), giving $d_1 = 19$ and
$d_2 = 217$. No other cascade quantity selects a dimension.
\end{proof}

\begin{center}
\begin{tabular}{ccrl}
\toprule
Index & $d$ & $\Omega_d$ & Role \\
\midrule
0 & 5 & 26.32 & Volume maximum (argmax $V_d$) \\
1 & 7 & 32.47 & Area maximum (natural zero of $p$) \\
2 & 19 & 0.516 & First threshold ($c_1 = \tfrac{1}{2}\ln\pi$) \\
3 & 217 & $2.335\times 10^{-120}$ &
Second threshold ($c_2 = \sqrt{\pi}$) \\
\bottomrule
\end{tabular}
\end{center}

\section{The Cascade Invariant}

\begin{theorem}[Content, not scale]\label{thm:content}
The natural combination of the threshold sphere areas is their product,
not their quotient.
\end{theorem}

\begin{proof}
By Corollary~\ref{cor:primary}, sphere area is the unique independent
cascade quantity, and any identification maps content to energy linearly:
$E_d = C\cdot\Omega_{d-1}$ with $C$ universal. Total content across two
independent levels is a product of two independent contributions. The
quotient measures relative scale, not total content.

The cascade originates at $d = \infty$ where
$V_\infty = \Omega_\infty = 0$: geometric nothing. As
$d_2\to\infty$: the product $\Omega_{d_1}\times\Omega_{d_2}\to 0$
(content approaches nothing---correct). The quotient
$\Omega_{d_1}/\Omega_{d_2}\to\infty$ (scale diverges at the origin---wrong).

The recurrence is first-order multiplicative
(Corollary~\ref{cor:primary}). Independent contributions at independent
levels combine as a product.
\end{proof}

\begin{definition}[Cascade invariant]
$I_0 := \Omega_{d_1}\times\Omega_{d_2} = \Omega_{19}\times\Omega_{217}$.
\end{definition}

\begin{theorem}[Uniqueness]\label{thm:unique}
$I_0$ is the unique scalar built from the threshold sphere areas under
three deductive constraints from the recurrence's structure.
\end{theorem}

\begin{proof}
(1)~\textit{Sphere areas, exponent~1.} The recurrence is first-order
multiplicative (Corollary~\ref{cor:primary}), outputting one copy of
$\Omega_d$ at each step. Higher powers $\Omega^k$ ($k\neq 1$) require
an operation not present in the recurrence.

(2)~\textit{Two factors.} Two thresholds
(Theorem~\ref{thm:thresholds}): two $\Omega$ factors.

(3)~\textit{Product.} The recurrence is multiplicative, so the natural
combination of its outputs is their product. The product measures content;
the quotient measures relative scale. Only the product has the correct
limiting behaviour at the cascade's origin
(Theorem~\ref{thm:content}).

The unique combination: $I_0 = \Omega_{d_1}\times\Omega_{d_2}$.
\end{proof}

\begin{theorem}[Evaluation]
$I_0 = 1.2051\times 10^{-120}$.
\end{theorem}

\begin{proof}
$\Omega_{19} = 0.51614$.\; $\Omega_{217} = 2.33490\times 10^{-120}$.\;
Product: $1.2051\times 10^{-120}$.
\end{proof}

\section{The Hierarchy}

\begin{theorem}\label{thm:hierarchy}
$\log_{10}(1/\Omega_{d_2})\approx
\pi e^{2\sqrt{\pi}}(2\sqrt{\pi}-1)/\ln 10\approx 120.3$.
\end{theorem}

\begin{proof}
Stirling at $d_2 = 2\pi e^{2\sqrt{\pi}}$:
$\ln\Omega_{d_2}\approx -\pi e^{2\sqrt{\pi}}(2\sqrt{\pi}-1)
\approx -277$ nats $\approx -120.3$ decades.
\end{proof}

\begin{remark}
$\sqrt{\pi}\approx 1.77$ (orthogonal compression);\;
$e^{2\sqrt{\pi}}\approx 34.6$ (exponential amplification);\;
$2\pi e^{2\sqrt{\pi}}\approx 218$ (the cascade dimension);\;
$(2\sqrt{\pi}-1)\approx 2.54$ (per-dimension suppression).\;
Product: $\approx 277$ nats $\approx 120$ decades.
\end{remark}

\section{Comparison with a Physical Constant}

The cascade invariant $I_0 = 1.2051\times 10^{-120}$ is a pure number from
the Gamma function. It was not constructed to match any physical quantity.

The observed ratio of the vacuum energy density to the fourth power of the
Planck mass is
$\rho_\Lambda/M^4_{\mathrm{Pl}} = (1.10\pm 0.02)\times 10^{-120}$,
determined by general relativity, quantum field theory, and astronomical
measurements---none of which enter this paper.

The agreement is $\sim\!10\%$, within $5\sigma$.

We record this without explaining it.

\section{Robustness}

\begin{center}
\begin{tabular}{llcl}
\toprule
Primitive set & Thresholds & $\log_{10} I$ & Deviation \\
\midrule
$\{c_1, c_2\}$ (this paper) & 19, 217 & $-119.92$ & $\sim 10\%$ \\
$\{0, c_2\}$ & 7, 217 & $-118.12$ & Factor 68 \\
$\{c_1, 2c_1\}$ & 19, 61 & $-17.3$ & $10^{3}$ OOM \\
$\{2c_1, c_2\}$ & 61, 217 & $-136.6$ & 17 OOM low \\
$\{c_1/2, c_2\}$ & 36, 217 & $-126.4$ & 6 OOM \\
$\{c_1, \pi\}$ & 19, $>$3000 & $<-1000$ & Irrelevant \\
$\{c_1, c_2\}$, quotient & 19, 217 & $+119.6$ & Wrong sign \\
\bottomrule
\end{tabular}
\end{center}

Shifting either threshold by one integer gives deviations of
70--900\%. No alternative primitive set drawn from operations present
in the recurrence comes within an order of magnitude of $10^{-120}$.

\section{Summary}

\begin{center}
\begin{tabular}{cll}
\toprule
Step & Statement & Source \\
\midrule
1 & $V_\infty = \Omega_\infty = 0$ & Definition \\
2 & Orthogonal $\Rightarrow\sqrt{\pi}$, irreducible & Thms 2.1, 4.1 \\
3 & $\Omega_{d-1}$ unique independent quantity & Cor 3.2 \\
4 & $d_V = 5$: volume maximum & argmax $V_d$ \\
5 & $d_0 = 7$: area maximum, natural zero & $p' > 0$ \\
6 & $c_1 = \tfrac{1}{2}\ln\pi$; unique across four classes & Thms 6.1, 6.2 \\
7 & $c_2 = \sqrt{\pi}$: second canonical value & Thm 6.3 \\
8 & Orbit terminates (first-order $\Rightarrow$ two forms) & Thm 6.5 \\
9 & $d_1 = 19$,\; $d_2 = 217$ & Thm 6.7 \\
10 & Four distinguished dimensions; no fifth & Thm 7.1 \\
11 & Content, not scale $\Rightarrow$ product & Thm 8.1 \\
12 & $I_0 = \Omega_{19}\times\Omega_{217}$; unique & Thm 8.3 \\
13 & $I_0 = 1.2051\times 10^{-120}$ & Computation \\
14 & $120\approx\pi e^{2\sqrt{\pi}}(2\sqrt{\pi}-1)/\ln 10$ & Thm 9.1 \\
\bottomrule
\end{tabular}
\end{center}

Every step is a theorem about the Gamma function. The derivation has zero
free parameters and no non-deductive steps.

\section{What This Paper Does Not Do}

\begin{itemize}
\item Explain the agreement with $\rho_\Lambda/M^4_{\mathrm{Pl}}$.
\item Introduce an observer, a hypothesis, or a physical interpretation.
\item Use any result from physics.
\end{itemize}

The paper is a theorem about the Gamma function.

\section*{Appendix: Numerical Verification}

$\sqrt{\pi} = 1.7724538509$\;;\quad
$\tfrac{1}{2}\ln\pi = 0.5723649429$

$\Omega_5 = \pi^3 = 26.319$\;;\quad
$\Omega_7 = \pi^4/3 = 32.470$\;;\quad
$\Omega_{19} = 0.51614$\;;\quad
$\Omega_{217} = 2.33490\times 10^{-120}$

$p(19) = 0.5535 < 0.5724 < p(20) = 0.6013$\;$\checkmark$

$p(217) = 1.77101 < 1.77245 < p(218) = 1.77331$\;$\checkmark$

$I_0 = 0.51614\times 2.33490\times 10^{-120}
= 1.2051\times 10^{-120}$\;$\checkmark$

$\pi e^{2\sqrt{\pi}}(2\sqrt{\pi}-1)/\ln 10 = 120.27$\;$\checkmark$

$R_{\mathrm{eff}}(5) = 1/\sqrt{8} = 0.354$\;;\quad
$R_{\mathrm{eff}}(19) = 1/\sqrt{22} = 0.213$\;;\quad
$R_{\mathrm{eff}}(217) = 1/\sqrt{220} = 0.067$

$\Omega_4/V_5 = 5$\;$\checkmark$\;;\quad
$\Omega_6/V_7 = 7$\;$\checkmark$

\end{document}